%=====================================================================
%  Project_Defense.tex
%  A defense (justification) of the ALGORITHM CHOICES for the project:
%  "AI-Powered Autonomous Smart Agriculture using Swarm Intelligence
%   and Reinforcement Learning".
%
%  Phase 1 : why PSO (and a Genetic Algorithm) for task allocation?
%  Phase 2 : why Value Iteration and Q-Learning for navigation?
%
%  Compile with:  pdflatex Project_Defense.tex
%=====================================================================
\documentclass[11pt]{article}

\usepackage[a4paper,margin=1in]{geometry}
\usepackage{amsmath}
\usepackage{amssymb}
\usepackage{booktabs}
\usepackage{xcolor}
\usepackage{tcolorbox}
\usepackage{enumitem}
\usepackage{hyperref}
\hypersetup{colorlinks=true, linkcolor=blue!60!black, urlcolor=blue!60!black}

% Coloured callout boxes (same family as the companion guides).
\newtcolorbox{claim}[1][]{colback=blue!4!white, colframe=blue!45!black,
  fonttitle=\bfseries, title=Claim, #1}
\newtcolorbox{rebut}[1][]{colback=green!4!white, colframe=green!40!black,
  fonttitle=\bfseries, title=Defense, #1}
\newtcolorbox{qbox}[2][]{colback=orange!5!white, colframe=orange!60!black,
  fonttitle=\bfseries, title=Likely examiner question: #2, #1}

\title{\textbf{Design Justification and Defense}\\[4pt]
\large Why These Algorithms? A Defense of the Phase-1 and Phase-2\\
Choices in the Smart-Agriculture Project}
\author{AI Laboratory Project}
\date{}

\begin{document}
\maketitle

\begin{abstract}
\noindent
This document \emph{defends} the algorithmic decisions taken in the project. It
argues that each algorithm was not chosen arbitrarily but follows directly from
the mathematical structure of the problem it solves. Phase~1 (task allocation)
is defended as a large, discrete, multi-objective combinatorial problem best
attacked by \textbf{population-based metaheuristics (PSO and a Genetic
Algorithm)}. Phase~2
(navigation) is defended as a \textbf{stochastic sequential decision problem
(an MDP)} best attacked by \textbf{reinforcement learning (Value Iteration and
Q-Learning)}. For each phase we state the problem's defining properties, justify
the chosen family, justify the specific algorithms, rule out the main
alternatives, and pre-empt the questions an examiner is most likely to ask.
\end{abstract}

\tableofcontents

%=====================================================================
\section{The Guiding Principle: Let the Problem Choose the Algorithm}
%=====================================================================

Our central defense is simple and consistent: \textbf{the structure of a problem
dictates which family of algorithms is appropriate}. We did not pick fashionable
algorithms; we identified two problems with fundamentally different structure and
applied the family that matches each.

\begin{center}
\renewcommand{\arraystretch}{1.4}
\begin{tabular}{@{}p{2.3cm} p{5.4cm} p{5.4cm}@{}}
\toprule
& \textbf{Phase 1 --- Allocation} & \textbf{Phase 2 --- Navigation} \\
\midrule
Question & \emph{Which} robot gets \emph{which} flowers? & \emph{How} does one robot move to its flowers? \\
Structure & Large discrete combinatorial optimization, multi-objective & Stochastic sequential decision process (MDP) \\
Search space & $5^{400}$ static assignments & Sequences of moves under uncertainty \\
Right family & \textbf{Population-based metaheuristics} & \textbf{Reinforcement learning / dynamic programming} \\
Our choice & PSO + Genetic Algorithm & Value Iteration + Q-Learning \\
\bottomrule
\end{tabular}
\end{center}

This clean separation --- allocation by population-based optimization, navigation
by reinforcement learning --- is itself a deliberate design decision, defended in
Section~\ref{sec:decomp}.

%=====================================================================
\section{Phase 1 Defense --- Why Population-Based Optimization for Allocation}
%=====================================================================

\subsection{The defining properties of the allocation problem}
Any defense of an algorithm must start from the problem. Ours has four
properties that together rule out most classical methods:

\begin{enumerate}[leftmargin=1.5em]
  \item \textbf{Combinatorially huge.} Assigning 400 flowers to 5 robots gives
        $5^{400}$ candidate solutions --- vastly more than the number of atoms in
        the universe. Exhaustive search is impossible.
  \item \textbf{Discrete.} A flower is assigned to robot 2 \emph{or} robot 3 ---
        there is no ``robot 2.5.'' Gradient-based optimizers (which need smooth,
        differentiable objectives) do not apply.
  \item \textbf{Multi-objective.} We simultaneously minimise travel distance,
        energy, completion time and workload imbalance while maximising coverage
        and robot utilisation --- six competing goals (Eq.~1 of the companion
        guide).
  \item \textbf{Non-linear \& coupled.} The cost of one robot's assignment
        depends on \emph{all} the flowers it received (a routing/TSP-like
        sub-problem), so the objective is highly non-linear and the flowers are
        not independent.
\end{enumerate}

\begin{claim}
A problem that is huge, discrete, multi-objective and non-linear is the textbook
definition of an \textbf{NP-hard combinatorial optimization problem}. This is
precisely the class for which \emph{population-based metaheuristics} such as
swarm intelligence and evolutionary algorithms were invented.
\end{claim}

\subsection{Why population-based metaheuristics specifically}
\begin{rebut}
\textbf{Metaheuristics} do not require derivatives, handle discrete variables
naturally, and find high-quality (not necessarily perfect) solutions in
reasonable time --- exactly what a $5^{400}$ space demands. Among them,
\textbf{population-based} methods (PSO, GA) are especially suitable because:
\begin{itemize}[nosep]
  \item they keep \textbf{many candidate allocations at once}, exploring in
        parallel and resisting getting stuck in one poor solution;
  \item they need only a \textbf{fitness value}, so our custom six-objective cost
        function plugs straight in with no reformulation;
  \item they are \textbf{robust and simple}, matching the ``from scratch, NumPy
        only'' constraint and the project theme of many simple agents/individuals
        producing an intelligent global result.
\end{itemize}
We deliberately chose \emph{two different families} of population method --- a
\textbf{swarm-intelligence} optimizer (PSO) and an \textbf{evolutionary}
optimizer (GA) --- so the comparison spans genuinely different search
mechanisms, not two variants of the same idea.
\end{rebut}

\subsection{Why Particle Swarm Optimization (PSO)}
\begin{rebut}
PSO is the natural \emph{first} choice because it is the canonical, well-studied
swarm optimizer with only three intuitive parameters ($w, c_1, c_2$). It treats
each candidate allocation as a whole and moves it through the search space using
personal-best and global-best memory. It is simple to implement from scratch,
fast, and gives a clear, reproducible convergence curve --- an ideal
\emph{baseline} optimizer for the allocation task.
\end{rebut}

\subsection{Why we \emph{also} added a Genetic Algorithm (GA)}
This is the strongest part of the Phase-1 defense, because it shows scientific
honesty rather than a single lucky choice.

\begin{claim}
Continuous PSO has a known weakness in \textbf{very high dimensions} (here, 400).
It encodes each flower as an independent continuous number and has \emph{no
built-in notion} that spatially close flowers should share a robot. In our
experiments PSO's allocation was only modestly better than random.
\end{claim}

\begin{rebut}
Rather than hide this, we introduced a \textbf{Genetic Algorithm as a second,
structurally different allocator}. GA operates \emph{directly} on the discrete
robot indices (no continuous-to-integer floor step), evolving a population via
\textbf{tournament selection}, \textbf{uniform crossover}, \textbf{mutation} and
\textbf{elitism}. We give it a strong start by \emph{seeding} part of the initial
population with a distance-greedy assignment (each flower to its nearest robot
depot), which injects spatial structure that selection then balances under the
fitness function. On the identical environment and identical fitness function,
GA reached a \emph{lower (better) fitness} ($1.43$ vs PSO's $1.59$), a shorter
total distance ($2204$ vs $2450$\,m) and a far more even workload (load std
$1.1$ vs $2.1$). Presenting both, and letting the data decide, is exactly how
algorithm selection should be defended: we \textbf{empirically justified} the
choice rather than merely asserting it.
\end{rebut}

\subsection{Ruling out the alternatives}

\begin{center}
\renewcommand{\arraystretch}{1.35}
\begin{tabular}{@{}p{3.6cm} p{9.4cm}@{}}
\toprule
\textbf{Alternative} & \textbf{Why we did not use it} \\
\midrule
Brute force / exhaustive & $5^{400}$ combinations --- computationally impossible. \\
Greedy (nearest robot) & Fast but myopic: overloads robots near dense clusters, ignores balance and global cost. We keep it as a \emph{baseline} to beat, and reuse it only to \emph{seed} the GA. \\
Integer Linear Programming & Would need a linear objective; ours is non-linear (routing inside each robot) and multi-objective. Also scales poorly and would violate the ``from scratch, NumPy only'' constraint. \\
Gradient descent & Requires a differentiable, continuous objective; assignment is discrete. \\
Ant Colony Optimization & Another valid metaheuristic, but its natural home is \emph{path/route} construction; for a pure assignment it adds pheromone machinery without a clear advantage over an evolutionary search. GA gives a cleaner, discrete comparison against PSO. \\
Deep learning & No labelled ``correct allocation'' dataset exists; this is an optimization problem, not a supervised-learning problem. \\
\bottomrule
\end{tabular}
\end{center}

\subsection{Anticipated examiner questions (Phase 1)}

\begin{qbox}{``Isn't PSO for continuous problems? Why use it on a discrete one?''}
Yes --- and we address this directly. We use the \textbf{standard continuous
encoding}: a particle is a real vector, and each component is floored to a robot
index. This is a well-established technique for applying PSO to discrete
assignment problems. We also \emph{acknowledge} its high-dimensional weakness and
back it up with a GA, which is discrete by construction. So the discreteness
concern is exactly \emph{why} the project includes two allocators.
\end{qbox}

\begin{qbox}{``If the GA wins, why keep PSO at all?''}
Because the project's value is the \textbf{comparison}. PSO is the canonical
swarm baseline; the GA is the evolutionary challenger. Showing \emph{when} and
\emph{why} one beats the other (continuous high-dimensional search vs.\ discrete
evolution with heuristic seeding) is a genuine research contribution and
demonstrates understanding of both algorithms' mechanics, not just their outputs.
\end{qbox}

\begin{qbox}{``Isn't seeding the GA with a greedy solution cheating?''}
No --- \textbf{heuristic initialisation} is a standard, textbook GA technique. The
greedy seed only provides a \emph{starting point}; it is unbalanced on its own
(one robot near a dense cluster is overloaded). The GA's selection, crossover and
mutation are what \emph{balance} the loads and lower the fitness. For fairness the
same greedy assignment is also available as an explicit baseline, and PSO is run
on the identical environment. The GA's win comes from evolution, not from the
seed alone.
\end{qbox}

\begin{qbox}{``How do you know your result is real and not luck?''}
Every experiment uses a \textbf{fixed random seed} and a \textbf{deterministic
greenhouse}, so results are byte-for-byte reproducible across runs. Both
allocators are judged by the \emph{same} fitness function on the \emph{same}
environment. The convergence curves and hyperparameter studies show the results
are stable, not cherry-picked.
\end{qbox}

%=====================================================================
\section{Phase 2 Defense --- Why Reinforcement Learning for Navigation}
%=====================================================================

\subsection{The defining properties of the navigation problem}
Phase~2 has a completely different structure, which is why it needs a completely
different algorithm family:

\begin{enumerate}[leftmargin=1.5em]
  \item \textbf{Sequential.} The robot makes a \emph{sequence} of dependent
        decisions (each move changes where the next move starts), not one static
        choice.
  \item \textbf{Stochastic.} Motors slip: an intended move succeeds only 80\% of
        the time. The outcome of an action is uncertain.
  \item \textbf{Delayed reward.} The big reward (reaching a flower) may come many
        steps after the moves that made it possible --- the algorithm must
        propagate credit backwards.
  \item \textbf{Goal-directed under constraints.} Avoid obstacles, conserve
        battery, and return to a charger.
\end{enumerate}

\begin{claim}
``Sequential decisions'' $+$ ``uncertain outcomes'' $+$ ``delayed reward'' is the
exact definition of a \textbf{Markov Decision Process}. The mathematically
correct tools for an MDP are \textbf{dynamic programming (Value Iteration)} and
\textbf{reinforcement learning (Q-Learning)} --- not a plain path-finder.
\end{claim}

\subsection{Why not just use A* or Dijkstra?}
This is the single most important question in the Phase-2 defense, so we answer
it head-on.

\begin{rebut}
A* and Dijkstra find the shortest path in a \textbf{deterministic} graph where
every move does exactly what you command. Our world is \textbf{stochastic} --- the
robot slips 20\% of the time. A single fixed A* path is brittle: one slip and the
robot is off the planned route with no instructions for the new cell. Value
Iteration and Q-Learning instead produce a \textbf{policy}: an optimal action for
\emph{every} cell, so the robot always knows what to do \emph{even after it
slips}. Handling uncertainty with a policy --- not a fixed path --- is precisely
why an MDP method is required rather than classical shortest-path search.
\end{rebut}

\subsection{Why Value Iteration (model-based)}
\begin{rebut}
When the environment is \textbf{known}, Value Iteration is the provably optimal,
efficient solution: it applies the Bellman optimality update until the values
converge and then extracts the optimal policy. It gives us a \textbf{gold-standard
benchmark} --- the best any method could possibly do on this MDP --- computed
quickly (\,$\sim$33 iterations\,) and reproducibly. Any learning method can then
be measured against it.
\end{rebut}

\subsection{Why Q-Learning (model-free)}
\begin{rebut}
Value Iteration's strength is also its limitation: it \emph{requires} a known
transition model. A real pollination robot dropped into a new greenhouse does
\textbf{not} know the obstacle map or its own slip probabilities in advance.
\textbf{Q-Learning} removes that assumption: it learns an optimal policy purely
from trial-and-error interaction, with no model at all. Including it demonstrates
the realistic, deployable case and lets us showcase the fundamental
\textbf{model-based vs model-free} trade-off --- the conceptual heart of
reinforcement learning.
\end{rebut}

\subsection{Why the two together are the right pairing}
\begin{claim}
Value Iteration and Q-Learning are the ``matched pair'' that isolates exactly one
variable: \emph{knowledge of the model}. Same grid, same obstacles, same
flowers, same rewards --- the \emph{only} difference is whether the model is
known.
\end{claim}
\begin{rebut}
This makes the comparison scientifically clean. We can show that Q-Learning,
starting from zero knowledge, \emph{rediscovers} nearly the same quality of route
that Value Iteration computes instantly --- validating both, and teaching the
core lesson of RL. That is a far stronger contribution than implementing either
one alone.
\end{rebut}

\subsection{Ruling out the alternatives}

\begin{center}
\renewcommand{\arraystretch}{1.35}
\begin{tabular}{@{}p{3.6cm} p{9.4cm}@{}}
\toprule
\textbf{Alternative} & \textbf{Why we did not use it} \\
\midrule
A* / Dijkstra / BFS & Assume deterministic moves; cannot handle 20\% slip. Produce a fragile fixed path, not a policy over all states. \\
Policy Iteration & A valid DP alternative to Value Iteration; converges to the same optimum. VI was chosen for its simpler single-loop update and cleaner convergence curve. (Equivalent, not inferior --- an acceptable substitute.) \\
Deep Q-Network (DQN) & Neural function approximation is unnecessary for a $20\times20$ grid --- a plain Q-table fits in memory and is exact. DQN would also violate the ``NumPy only, no TensorFlow/PyTorch'' constraint and add non-reproducible training noise. \\
SARSA & On-policy cousin of Q-Learning; reasonable, but Q-Learning's off-policy update learns the optimal greedy policy more directly and pairs more cleanly with Value Iteration's optimality. \\
Supervised learning & There is no dataset of ``correct moves'' to imitate; the robot must \emph{discover} good behaviour from rewards. \\
\bottomrule
\end{tabular}
\end{center}

\subsection{Anticipated examiner questions (Phase 2)}

\begin{qbox}{``This is just a maze --- why not A*? It's simpler.''}
A* solves a \emph{deterministic} maze. Ours is stochastic (Section 4.2). A* gives
one path; a slip breaks it. We need an action for every possible cell the robot
could slip into, i.e.\ a \emph{policy}, which is what Value Iteration and
Q-Learning produce. The stochastic transition model is the reason a plain maze
solver is insufficient.
\end{qbox}

\begin{qbox}{``Why does Q-Learning need a bitmask? A* wouldn't.''}
Because the robot must visit \emph{several} flowers and an MDP state is
memoryless. Without a record of which flowers are already done, the robot cannot
know when the mission is complete or avoid re-visiting. The 6-bit visited-mask
augments the state with exactly that memory, letting a memoryless agent complete
a multi-goal mission. It is a standard and necessary technique, not a hack.
\end{qbox}

\begin{qbox}{``Q-Learning is slower than Value Iteration --- isn't it worse?''}
No --- they optimise different assumptions. VI is faster \emph{because} it is
handed the full model; Q-Learning is slower \emph{because} it must first learn
the model implicitly from experience. In a real unknown greenhouse, VI cannot run
at all (no model), while Q-Learning can. The ``slower'' cost buys
\textbf{model-freedom}, which is the realistic deployment scenario.
\end{qbox}

%=====================================================================
\section{Defending the Two-Phase Decomposition Itself}
\label{sec:decomp}
%=====================================================================

An examiner may ask why we split the problem at all rather than solving
everything with one algorithm.

\begin{rebut}
Allocation and navigation are \textbf{structurally different problems}
(Section 1): one is a static combinatorial assignment, the other a dynamic
stochastic control problem. Forcing a single algorithm to do both would mean
using the wrong tool for one of them. \textbf{Decomposition} is a well-established
engineering principle: solve each sub-problem with the method that fits it, then
compose. Phase~1 outputs the flower assignments; Phase~2 consumes them as fixed
input. The interfaces are clean (a CSV of assignments), the algorithms never mix,
and each phase can be validated independently. This modularity also makes the
whole pipeline reproducible and easy to reason about.
\end{rebut}

%=====================================================================
\section{Summary of the Defense}
%=====================================================================

\begin{center}
\renewcommand{\arraystretch}{1.4}
\begin{tabular}{@{}p{2.3cm} p{4.6cm} p{5.8cm}@{}}
\toprule
\textbf{Phase} & \textbf{Chosen algorithms} & \textbf{One-line justification} \\
\midrule
Phase 1 & PSO + Genetic Algorithm & The allocation is a huge, discrete, multi-objective combinatorial problem --- the home ground of population-based optimization. PSO is the canonical swarm baseline; the GA's discrete encoding plus heuristic seeding fixes PSO's high-dimensional weakness and wins empirically. \\
Phase 2 & Value Iteration + Q-Learning (RL) & Navigation is a stochastic MDP with delayed reward, so it needs a \emph{policy}, not a fixed path. VI is the model-based gold standard; Q-Learning is the model-free realistic case; together they isolate the model-based vs model-free trade-off. \\
\bottomrule
\end{tabular}
\end{center}

\begin{tcolorbox}[colback=green!5!white,colframe=green!45!black,title=Closing statement]
Every algorithm in this project is chosen because the \textbf{structure of the
problem demands it}, and every choice is backed by a \textbf{fair, reproducible,
data-driven comparison} against a meaningful alternative. Phase~1 defends
population-based optimization for allocation and empirically prefers the Genetic
Algorithm over PSO; Phase~2 defends reinforcement learning for navigation and
demonstrates both the model-based (Value Iteration) and model-free (Q-Learning)
sides of that family. The two-phase decomposition ensures each problem is solved
by the algorithm best suited to it.
\end{tcolorbox}

\end{document}
